\usepackage{geometry}
\usepackage{fontspec}
\usepackage{amsmath, amsthm, amssymb}
\usepackage{mathrsfs}
\usepackage{enumitem}
\setlist{nosep}
\usepackage{graphicx}
\usepackage{array}
\usepackage{ulem}
\usepackage{caption}
\usepackage{tocloft}

% 2. 图形和颜色宏包
\usepackage[dvipsnames]{xcolor}
\usepackage{tikz}
\usetikzlibrary{shapes.geometric,arrows.meta,calc,intersections,positioning}
\usepackage[most]{tcolorbox}
\tcbuselibrary{theorems}

% 3. 页眉页脚宏包
\usepackage{fancyhdr}
\usepackage{lastpage}

% 4. 超链接和书签
\usepackage{hyperref}
\usepackage{bookmark}

% 5. 智能引用
\usepackage{cleveref}

\usepackage{draftwatermark}

% 设置水印文本
\SetWatermarkText{还在尬黑}
% 设置水印缩放比例
\SetWatermarkScale{1}
% 设置水印颜色和透明度
\SetWatermarkColor[gray]{1}
% 设置水印角度
\SetWatermarkAngle{45}

% 1. 页面布局设置（窄边距）
\geometry{
  a4paper,
  top=15mm,
  bottom=10mm,
  left=15mm,
  right=15mm,
  headheight=25pt,
  headsep=8mm,
  footskip=15mm,
  includehead,
  includefoot
}

% 2. 字体设置
\setmainfont{Times New Roman}
\setsansfont{Arial}
\setmonofont{Consolas}

% 定义字体命令
\newcommand{\kt}{\kaishu}
\newcommand{\st}{\songti}
\newcommand{\dd}{\text{d}}

\definecolor{deepblue}{HTML}{003366}
\definecolor{lightblueA}{HTML}{E5F6FF}
\definecolor{deepgreen}{HTML}{006633}
\definecolor{lightgreen}{HTML}{E5F6E5}
\definecolor{deeppurple}{HTML}{660066}
\definecolor{lightpurple}{HTML}{F6E5F6}
\definecolor{deepred}{HTML}{990000}
\definecolor{lightred}{HTML}{FFE5E5}
\definecolor{deeporange}{HTML}{CC6600}
\definecolor{lightorange}{HTML}{FFF6E5}
\definecolor{deepgray}{HTML}{333333}
\definecolor{lightgrayA}{HTML}{F6F6F6}
\definecolor{deepbrown}{HTML}{8B4513}
\definecolor{lightbrown}{HTML}{FFE4B5}
\definecolor{lightyellow}{HTML}{FFFFE0}
\definecolor{darkyellow}{HTML}{CCCC00}
\definecolor{lightcyan}{HTML}{E0FFFF}
\definecolor{darkcyan}{HTML}{008B8B}
\definecolor{lightpink}{HTML}{FFB6C1}
\definecolor{darkpink}{HTML}{FF1493}
\definecolor{lavender}{HTML}{E6E6FA}
\definecolor{thistle}{HTML}{D8BFD8}
\definecolor{lightblueB}{HTML}{ADD8E6}
\definecolor{darkblue}{HTML}{00008B}
\definecolor{lightgrayB}{HTML}{D3D3D3}
\definecolor{darkgray}{HTML}{A9A9A9}

\newcounter{theoremcounter}[section]

% 定理环境
\newtcbtheorem[use counter=theoremcounter,number within=section]{theorem}{定理}{
    colback=lightgreen,
    colframe=deepgreen,
    colbacktitle=deepgreen,
    coltitle=white,
    fonttitle=\upshape\color{white},
    fontupper=\rmfamily,
    separator sign none,
    top=8pt,
    attach boxed title to top left={yshift=-2mm,xshift=5mm},
    description delimiters={ }{ },
    enhanced
}{thm}

% 引理环境
\newtcbtheorem[use counter=theoremcounter,number within=section]{lemma}{引理}{
    colback=lightpurple,
    colframe=deeppurple,
    colbacktitle=deeppurple,
    coltitle=white,
    fonttitle=\upshape\color{white},
    fontupper=\rmfamily,
    separator sign none,
    top=8pt,
    attach boxed title to top left={yshift=-2mm,xshift=5mm},
    description delimiters={ }{ },
    enhanced
}{lem}

% 定义环境
\newtcbtheorem[number within=section]{definition}{定义}{
    colback=lightred,
    colframe=deepred,
    colbacktitle=deepred,
    coltitle=white,
    fonttitle=\upshape\color{white},
    fontupper=\rmfamily,
    separator sign none,
    top=8pt,
    attach boxed title to top left={yshift=-2mm,xshift=5mm},
    description delimiters={ }{ },
    enhanced
}{def}

% 例题和解共享计数器
\newcounter{examplecounter}[section]

% 例题环境
\newtcbtheorem[use counter=examplecounter,number within=section]{example}{例题}{
    colback=lightblueA,
    colframe=deepblue,
    colbacktitle=deepblue,
    coltitle=white,
    fonttitle=\upshape\color{white},
    fontupper=\rmfamily,
    separator sign none,
    top=8pt,
    attach boxed title to top left={yshift=-2mm,xshift=5mm},
    description delimiters={ }{ },
    enhanced
}{ex}

% 结论环境
\newtcbtheorem[number within=section]{conclusion}{结论}{
    colback=lightorange,
    colframe=deeporange,
    colbacktitle=deeporange,
    coltitle=white,
    fonttitle=\upshape\color{white},
    fontupper=\rmfamily,
    separator sign none,
    top=8pt,
    attach boxed title to top left={yshift=-2mm,xshift=5mm},
    description delimiters={ }{ },
    enhanced
}{con}

\newtcbtheorem[number within=section]{property}{性质}{
    colback=lightgrayA,
    colframe=deepgray,
    colbacktitle=deepgray,
    coltitle=white,
    fonttitle=\upshape\color{white},
    fontupper=\rmfamily,
    separator sign none,
    top=8pt,
    attach boxed title to top left={yshift=-2mm,xshift=5mm},
    description delimiters={ }{ },
    enhanced
}{prop}

\newtcbtheorem[number within=section]{criterion}{准则}{
    colback=lightbrown,
    colframe=deepbrown,
    colbacktitle=deepbrown,
    coltitle=white,
    fonttitle=\upshape\color{white},
    fontupper=\rmfamily,
    separator sign none,
    top=8pt,
    attach boxed title to top left={yshift=-2mm,xshift=5mm},
    description delimiters={ }{ },
    enhanced
}{crit}

\newtcbtheorem[use counter=theoremcounter,number within=section]{corollary}{推论}{
    colback=lightcyan,
    colframe=darkcyan,
    colbacktitle=darkcyan,
    coltitle=white,
    fonttitle=\upshape\color{white},
    fontupper=\rmfamily,
    separator sign none,
    top=8pt,
    attach boxed title to top left={yshift=-2mm,xshift=5mm},
    description delimiters={ }{ },
    enhanced
}{cor}

% 传统样式环境（无背景色）
\newtheoremstyle{plain-chinese}
  {6pt}
  {6pt}
  {\st}
  {}
  {\heiti}
  {.}
  { }
  {}
\theoremstyle{plain-chinese}
\newtheorem{solution}{解}[section]
\newtheorem{remark}{注}[section]
\newtheorem{myproof}{证明}[section]

% 智能引用设置
\crefname{theorem}{定理}{定理}
\crefname{lemma}{引理}{引理}
\crefname{definition}{定义}{定义}
\crefname{example}{例题}{例题}
\crefname{conclusion}{结论}{结论}
\crefname{solution}{解}{解}
\crefname{remark}{注}{注}
\crefname{property}{性质}{性质}
\crefname{criterion}{准则}{准则}
\crefname{myproof}{证明}{证明}
\crefname{corollary}{推论}{推论}

% 数学字体设置
\usepackage{unicode-math}
\setmathfont{Latin Modern Math}

% ========== 自定义命令 ==========
\newcommand{\R}{\mathbb{R}}
\newcommand{\C}{\mathbb{C}}
\newcommand{\Z}{\mathbb{Z}}
\newcommand{\N}{\mathbb{N}}
\newcommand{\dif}{\mathrm{d}}
\newcommand{\PP}{\mathbb{P}}
\newcommand{\RP}{\mathbb{RP}}
\newcommand{\CP}{\mathbb{CP}}
\newcommand{\A}{\mathbb{A}}
\newcommand{\F}{\mathbb{F}}
\newcommand{\Span}{\mathrm{span}}
\newcommand{\rank}{\mathrm{rank}}
\newcommand{\diag}{\mathrm{diag}}

% ========== 图形路径设置 ==========
\graphicspath{{./flg/}}

% 页眉页脚设置
\pagestyle{fancy}
\fancyhf{}

\fancyhead[L]{\small\kt 高中数学~解析几何}
\fancyhead[C]{\small\st 居敬持志~守正出奇}
\fancyhead[R]{\small\st 还在尬黑出品}
\fancyfoot[C]{\thepage}

\fancyhead[R]{\small\st\rightmark}

\fancypagestyle{frontmatter}{
    \fancyhf{}
    \fancyhead[L]{\small\kt 高中数学~解析几何}
    \fancyhead[C]{\small\st 前言}
    \fancyhead[R]{\small\st 还在尬黑出品}
    \fancyfoot[C]{\thepage}
    \renewcommand{\headrulewidth}{0.4pt}
    \renewcommand{\footrulewidth}{0pt}
}

\fancypagestyle{tocmatter}{
    \fancyhf{}
    \fancyhead[L]{\small\kt 高中数学~解析几何}
    \fancyhead[C]{\small\st 目录}
    \fancyhead[R]{\small\st 还在尬黑出品}
    \fancyfoot[C]{\thepage}
    \renewcommand{\headrulewidth}{0.4pt}
    \renewcommand{\footrulewidth}{0pt}
}

\fancypagestyle{mainmatter}{
    \fancyhf{}
    \fancyhead[L]{\small\kt 高中数学~解析几何}
    \fancyhead[C]{\small\st 居敬持志~~守正出奇}
    \fancyhead[R]{\small\st\rightmark}
    \fancyfoot[C]{\thepage}
    \renewcommand{\headrulewidth}{0.4pt}
    \renewcommand{\footrulewidth}{0pt}
}

\renewcommand{\chaptermark}[1]{\markboth{#1}{}}
\renewcommand{\sectionmark}[1]{\markright{\thesection.\ #1}}

% 章节标题字体设置
\ctexset{
    chapter = {
        format = \centering,
        nameformat = \kaishu\LARGE,
        titleformat = \songti\LARGE,
        aftername = \quad,
        beforeskip = 30pt,
        afterskip = 20pt,
        name = {第,章},
        number = \chinese{chapter},
    },
    section = {
        format = \raggedright,
        nameformat = \heiti\large,
        titleformat = \songti\large,
        aftername = \quad,
        beforeskip = 15pt,
        afterskip = 6pt,
    },
    subsection = {
        format = \raggedright,
        nameformat = \heiti\normalsize,
        titleformat = \songti\normalsize,
        aftername = \quad,
        beforeskip = 6pt,
        afterskip = 3pt,
    }
}